%% kapa: exam=1 round=31 year=2020 period=2 form=A subject=appraisal_law subject_ko=감정평가관계법규 seq=26 number=26
%% source: Q-Net 감정평가사 기출문제 (cst00309) artlSeq=5207219 file=extracted/1008/5207219/A형/A형/2020년도 제31회 감정평가사 1차 시험_2교시_A형.pdf
%% conversion: kordoc 4.24.51 -> tools/gen1.py (auto); review=proofread-sonnet; flags=none
\begin{question}{26}
건축법령상 면적이 200제곱미터 이상인 대지에 건축을 하는 건축주는 용도지역 및 건축물의 규모에 따라 해당 지방자치단체의 조례로 정하는 기준에 따라 대지에 조경이나 그 밖에 필요한 조치를 하여야 한다. 다만, 건축법령은 예외적으로 조경 등의 조치를 필요로 하지 않는 건축물을 허용하고 있다. 이러한 예외에 해당하는 것을 모두 고른 것은? (단, 그 밖의 조례, 「건축법」 제73조에 따른 적용 특례, 건축협정은 고려하지 않음)
\begin{framedtext}
ㄱ. 축사 \par
ㄴ. 녹지지역에 건축하는 건축물 \par
ㄷ. 「건축법」상 가설건축물 \par
ㄹ. 면적 4천 제곱미터인 대지에 건축하는 공장 \par
ㅁ. 상업지역에 건축하는 연면적 합계가 1천 500제곱미터인 물류시설 \par
\end{framedtext}
\begin{choices}
  \item ㄱ, ㄴ, ㄹ
  \item ㄱ, ㄴ, ㅁ
  \item ㄷ, ㄹ, ㅁ
  \item ㄱ, ㄴ, ㄷ, ㄹ
  \item ㄴ, ㄷ, ㄹ, ㅁ
\end{choices}
\end{question}
\answer{4}
